\section{Fermat's Christmas theorem}
\lecture\label{lec:Fermat}

It is time to give a very nice number-theoretical application.

\begin{theorem}[Fermat]
\index{Fermat's Christmas theorem}
\label{thm:Fermat}
	Let $p\in\Z_{>0}$ be a prime number. The following statements are equivalent:
	\begin{enumerate}
	    \item $p=2$ or $p\equiv1\bmod 4$.
	    \item There exists $a\in\Z$ such that $a^2\equiv-1\bmod p$.
	    \item $p$ is not irreducible in $\Z[i]$.
	    \item $p=a^2+b^2$ for some $a,b\in\Z$.
	\end{enumerate}
\end{theorem}

\index{Fermat's little theorem}
\begin{proof}
    We first prove that $1)\implies 2)$. If $p=2$, take $a=1$. If $p=4k+1$ for
    some $k\geq1$, then by Fermat's little theorem, the polynomial 
    $X^{p-1}-1\in(\Z/p)[X]$ has roots $1,2,\dots,p-1$. Write
    \[
    (X-1)(X-2)\cdots (X-(p-1))=X^{p-1}-1=X^{4k}-1=(X^{2k}+1)(X^{2k}-1)
    \]
    in $(\Z/p)[X]$. Since $p$ is prime, $\Z/p$ is a field, and hence
    $(\Z/p)[X]$ is a unique factorization domain (because it is euclidean).
    Thus there exists $\alpha\in\Z/p$ such that $\alpha^{2k}+1=0$. To finish
    the proof choose an integer $a$ whose residue modulo $p$ is $\alpha^k$.
    Then $a^2\equiv-1\bmod p$.
    
    We now prove that $2)\implies 3)$. If $a^2\equiv-1\bmod p$, then $a^2+1=kp$ 
    for some $k\in\Z$. Since $(a-i)(a+i)=a^2+1=kp$, it follows that
    $p$ divides $(a-i)(a+i)$. Let us prove that $p$ is not prime in $\Z[i]$. 
    We claim that $p$ does not divide $a-i$ in $\Z[i]$. Indeed, if $p\mid a-i$, then
    $a-i=p(e+fi)$ for some $e,f\in\Z$ and this implies that $-1=pf$, a contradiction. Similarly,
    $p$ does not divide $a+i$. Thus $p$ is not prime in $\Z[i]$ 
    and hence it is not irreducible in $\Z[i]$ (because in $\Z[i]$ primes and irreducibles coincide). 
    
    We now prove that $3)\implies 4)$. Assume that $p=(a+bi)(c+di)$.  Since
    $a+bi\not\in\mathcal{U}(\Z[i])$ and $c+di\not\in\mathcal{U}(\Z[i])$, their
    norms are greater than one.  Since 
    \[
    p^2=(a^2+b^2)(c^2+d^2)
    \]
    in $\Z$, unique factorization in $\Z$ implies that $p=a^2+b^2$. 
    
    Finally, we prove that $4)\implies 1)$. 
    The only possible remainders after division by four are $0,1,2$ and $3$.  
    For all $a$, either $a^2\equiv 0\bmod 4$ or $a^2\equiv 1\bmod 4$. 
    If $p\equiv3\bmod 4$, then $p$ is never a sum of two squares, as 
    $a^2+b^2\equiv 0\bmod 4$, $a^2+b^2\equiv 1\bmod 4$ or $a^2+b^2\equiv 2\bmod 4$. 
    Therefore $p\not\equiv3\bmod4$. Since $p$ is prime, either
    $p=2$ or $p\equiv1\bmod4$.
\end{proof}

\begin{optional}
An alternative proof of the implication 
$1)\implies 2)$ goes as follows: if $p=4k+1$ is prime, then 
$x=(2k)!$ is such that $x^2\equiv -1\bmod p$. 
By 
Wilson's theorem, 
\[
(p-1)! = (4k)!\equiv -1\bmod p.
\]
Moreover,  
\begin{align*}
(4k)!&\equiv(4k)(4k-1)\cdots(2k+2)(2k+1)(2k)(2k-1)\cdots1\bmod p\\
&\equiv-1(-2)\cdots(-2k+1)(-2k)(2k)(2k-1)\cdots1\bmod p\\
&\equiv(-1)^{2k}(2k)!(2k)!\bmod p.
\end{align*}
%Hence $x=(2k)!$ is such that $x^2\equiv -1\bmod p$.
\end{optional}

\begin{exercise}
    Decompose the prime 41 as a sum of two squares. 
\end{exercise}

In 1625, Girard stated without proof a result characterizing which integers can
be expressed as the sum of two squares. Later, in a letter to Mersenne dated
December 25, 1640, Fermat presented several improvements and claimed to have proved
the result (this is the ``Christmas Letter''). Euler later published the first
rigorous proof, explicitly acknowledging Fermat's earlier work.

\begin{optional}
The previous exercise can be solved by using computers and a beautiful formula
of Gauss that uses binomial coefficients. Let $p=4k+1$ be a prime number.  If 
\[
a=\left\langle \frac12\binom{2k}{k}\right\rangle, 
\quad
b=\langle (2k)!a\rangle,
\]
where $\langle x\rangle$ denotes the 
residue of $x$ modulo $p$ of absolute value smaller than $p/2$, 
then it follows that $p=a^2+b^2$. Simple and beautiful as it is, Gauss' formula does not seem to be of any real computational value. 
\end{optional}

In \cite{MR1041889}, the author employs the Euclidean algorithm to determine efficiently 
the decomposition of a prime number $p=4k+1$ as the sum of two squares.

\begin{exercise}
\label{xca:p=3(4)}
    Let $\alpha\in\Z[i]$ be such that $\norm(\alpha)=p^2$ 
    for some prime number $p\in\Z$ with $p\equiv3\bmod 4$. 
    Prove that $\alpha$ is irreducible. 
\end{exercise}

\begin{exercise}
    \label{xca:p=1(4)}
    Let $p\in\Z$ be a prime number such that $p\equiv1\bmod 4$.
    Then $p=\alpha\overline{\alpha}$
    for some $\alpha\in\Z[i]$ irreducible. 
\end{exercise}

\begin{exercise}
\label{xca:Z[i]irreducibles}
Find the irreducible elements of $\Z[i]$. 
\end{exercise}

As a consequence of Theorem \ref{thm:Fermat} one can prove the following
result: An integer $n\geq2$ can be written as a sum of two squares
if and only if every prime $p\equiv3\bmod4$ occurs with even
exponent in the prime factorization of $n$.

\begin{optional}
The numbers that can be represented as the sums of two squares form the integer sequence \href{https://oeis.org/A001481}{A001481}: 
\[
0, 1, 2, 4, 5, 8, 9, 10, 13, 16, 17, 18, 20, 25, 26, 29, 32,\dots 
\] 
They form the set of all norms of Gaussian integers. 
\end{optional}

\begin{bonus}[Eisenstein integers]
\index{Eisenstein integers}
    %This exercise is about \emph{Eisenstein integers}. 
    \label{xca:Eisenstein}
    Let $\omega=e^{2\pi i/3}\in\C$ and 
    \[
    R=\Z[\omega]=\{a+b\omega:a,b\in\Z\}.
    \]
    Prove the following statements:
    \begin{enumerate}
        \item The map 
        \[
        N\colon R\to\Z_{\geq0},
        \quad
        N(a+b\omega)=a^2-ab+b^2,
        \]
        is multiplicative and 
        satisfies $N(\alpha)=\alpha\overline{\alpha}$
        for all $\alpha\in R$. 
        \item $\mathcal{U}(R)=\{1,-1,\omega,-\omega,\omega^2,-\omega^2\}$.
        \item $R$ is a euclidean domain.
        \item If $N(\alpha)$ is a prime number, then $\alpha$ is irreducible. 
        \item If $N(\alpha)=p^2$ for some 
            prime number $p$ with $p\equiv 2\bmod 3$, then 
            $\alpha$ is irreducible. 
        \item If $p$ is a prime number such that $p\equiv 1\bmod 3$, then 
            $p=\gamma\overline{\gamma}$ for some 
            irreducible $\gamma\in R$.
        \item Up to multiplication by units, 
            the irreducible elements of $R$ are $1-\omega$, 
            $a+b\omega$ with $a^2-ab+b^2=p$ for some prime number 
            $p\equiv 1\bmod 3$, 
            and prime numbers $p\in\Z$ with $p\equiv 2\bmod 3$.
    \end{enumerate}
\end{bonus}


